% Imporditud: tools/import_eksperimendid.py
% Allikas: Eesti füüsikaolümpiaadi eksperimendi ülesannete
%   kogu (Rael Kalda, 2023), ülesanne Ü112, lahendus L112.
\setAuthor{Jaan Kalda}
\setRound{lõppvoor}
\setYear{2015}
\setNumber{P E2}
\setDifficulty{3}
\setTopic{staatika}
\setVahendid{tasane plaat, joonlaud, tühi papist joogitops}
\setPunktid{12}
\setAllikas{Eksperimendi kogu Ü112, lahendus lk 88}
\setLahendusseis{toores}

\prob{Tops}
Määrata võimalikult täpselt joogitopsi massikeskme kaugus lauapinnast siis, kui tops seisab püsti laual.

\hint

\solu
Paneme topsi kaldpinnale ja suurendame kallet niikaua kuni tops ümber kukub. Mõõdame sel hetkel kaldpinna otste kõrguste vahe h (joonisel |EF | = h). Mõõdame topsi alumise serva diameetri d (täpsemalt: toetuspinnaga moodustuva kontaktjoone välise diameetri), joonisel d = 2|AO|. Tops kukub ümber, kui massikeskme vertikaalprojektsioon läheb kontaktjoonest väljapoole, piirjuhul on see otse kontaktjoone kohal, vt joonist. Seetõttu on joonisel $\angle$ACO = $\angle$EBF vastavalt ristuvate külgede tõttu ning seega on kolmnurgad ACO ja EBF sarnased. Seega otsitav kõrgus |OC| = |BF | $\cdot$ |AO|/|EF | = $L_2$ $-$ h2d/2h, kus L on laua pikkus |BE| (samuti mõõdame üle).

E

C

O A B

F
\probend
